\documentclass[10pt,a4paper]{amsart}
\usepackage[margin=19mm]{geometry}
\usepackage{amsmath,amssymb,mathtools}
\usepackage[scr=rsfs,cal=euler]{mathalfa}
\usepackage{xfrac}
\usepackage[unicode,hidelinks,bookmarks=false]{hyperref}
\newcommand{\ie}{i.e.\ }
\newcommand{\cf}{cf.\ }
\newcommand{\resp}{respectively\ }
\newcommand{\Subsection}{\subsection}
\providecommand{\nobreakdash}{\nobreak\hbox{-}}
\setlength{\emergencystretch}{2em}
\multlinegap0pt
\usepackage{bm}
\usepackage{bbm}
\usepackage{dsfont}
\usepackage{xspace}
\usepackage{cancel}
\usepackage{mathtools}
\usepackage{amsthm}
\makeatletter
\@ifundefined{theorem}{\newtheorem{theorem}{Theorem}}{}
\@ifundefined{proposition}{\newtheorem{proposition}[theorem]{Proposition}}{}
\makeatother
\newcommand{\E}{\mathbb{E}}
\newcommand{\R}{\mathbb{R}}
\newcommand{\Yt}{Y_t}
\newcommand{\bI}{\bm{I}}
\newcommand{\sYt}{s_{\Yt}}
\title[Information Gradient for Nonlinear Gaussian Channel with App]{Information Gradient for Nonlinear Gaussian Channel with Applications to Task-Oriented Communication}
\author{Tadashi Wadayama}
\date{Source: arXiv:2510.20179v1 (CC BY 4.0)}
\begin{document}
\maketitle

\noindent\emph{Source and licence.} Information Gradient for Nonlinear Gaussian Channel with Applications to Task-Oriented Communication. Version arXiv:2510.20179v1 (CC BY 4.0), distributed under the Creative Commons Attribution 4.0 International licence, as recorded in the arXiv licence metadata for this exact version and shown on the abstract page https://arxiv.org/abs/2510.20179v1. Full rights notice: licenses/arxiv-2510.20179-CC-BY-4.0.txt.\par\smallskip

\noindent\textit{Editorial scope.} Counted unit: the Proposition whose source label is \texttt{prop:info\_gradient} in the article (source0.tex lines 329--341, source label \texttt{prop:info\_gradient}), delivered together with the article's own complete original proof of it (source0.tex lines 343--400, one \texttt{proof} environment of 58 lines). No mathematics is written, repaired or re-derived: statement and proof are the article's own text line for line. Required context retained uncounted: none. Every definition, display and label of those lines is retained unchanged. Omitted from this package and not redistributed: 1--328, 342--342, 401--1845. Nothing inside a retained excerpt is omitted or altered: every retained body is delivered exactly as the article's lines are written, so no omission receipt of any kind accompanies this package. Citations: the retained excerpts carry no citation command at all, so no bibliography entry of the article is transcribed and no reference is redistributed. Reference adaptation: none was needed and none was made; every reference inside the delivered unit resolves inside it, so no unresolved reference or citation is produced. Printed slips kept literal: no printed word, symbol, hypothesis or number of the counted statement or of its proof was altered. Layout and class: the article's own class is \texttt{IEEEtran} with packages including amsmath, amssymb, amsfonts, amsthm, textcomp, xcolor, cite, bbm, bm, hyperref, graphicx, algorithm, algpseudocode; the wrapper is the lane's pinned \texttt{amsart} editorial wrapper at 10pt on A4, which restates the theorem-like environments the retained text needs and adds the article's own macro definitions that the retained text needs (\texttt{\textbackslash E}, \texttt{\textbackslash R}, \texttt{\textbackslash Yt}, \texttt{\textbackslash bI}, \texttt{\textbackslash sYt}), with extra wrapper packages (bm, bbm, dsfont, xspace, cancel, mathtools). The delivered numbering is the unit's own. Assets: no retained span contains a figure environment, a graphics-inclusion command, a PDF-page inclusion, a table or a TikZ picture; no image file is copied into this package, the package declares no asset dependency and no image file is redistributed with it. Licence: version arXiv:2510.20179v1 of this article is distributed under the Creative Commons Attribution 4.0 International licence, as recorded in the arXiv licence metadata for this exact version and shown on the abstract page https://arxiv.org/abs/2510.20179v1. A full rights notice is delivered at licenses/arxiv-2510.20179-CC-BY-4.0.txt. Characters: every retained excerpt is pure ASCII, so the delivered engine, XeLaTeX with its default Latin Modern text font, reports no missing character.
\begin{proposition}[Information gradient]
\label{prop:info_gradient}
Under standard regularity conditions, the gradient of mutual information 
with respect to the front-end parameters $\bm \eta$ is given by
\begin{align}
\nabla_{\bm \eta} I(X; \Yt) = -\E_{X,Z_t}\left[D f_{\bm \eta}(X)^\top \sYt(f_{\bm \eta}(X) + Z_t)\right],
\label{eq:gradient_formula}
\end{align}
where $\sYt(\bm y) \equiv \nabla_y \log p_{\Yt,\eta}(\bm y)$ 
is the score function of the marginal distribution of $\Yt$.
The notation $D f_{\bm \eta}(\bm x) \in \R^{m \times d}$ denotes 
the Jacobian of $f_{\bm \eta}$ with respect to $\bm \eta$.
\end{proposition}

\begin{proof}
Let $W = f_{\bm \eta}(X)$. Since $\Yt$ depends on $X$ only through $W$, the triple $X \to W \to \Yt$ forms a Markov chain. By the data processing inequality and the chain rule for mutual information:
\begin{align}
I(X; \Yt) = I(W; \Yt) + I(X; \Yt|W) = I(W; \Yt),
\label{eq:markov_chain}
\end{align}
where the second equality follows from the conditional independence $X \perp \Yt \mid W$.

The mutual information can be decomposed as
\begin{align}
I(W; \Yt) = h(\Yt) - h(\Yt|W).
\label{eq:mi_decomp_proof}
\end{align}
Given $W = \bm w$, the conditional distribution is $\Yt|W=\bm w \sim \mathcal{N}(\bm w, t\bI_m)$, 
hence we have
\begin{align}
h(\Yt|W=\bm{w}) = \frac{m}{2}\log(2\pi e t),
\end{align}
which is independent of both $\bm w$ and $\bm \eta$. Therefore:
\begin{align}
\nabla_{\bm \eta} h(\Yt|W) = \bm{0}.
\label{eq:cond_entropy_const}
\end{align}
Combining \eqref{eq:markov_chain}, \eqref{eq:mi_decomp_proof}, and \eqref{eq:cond_entropy_const},
we have
\begin{align}
\nabla_{\bm \eta} I(X; \Yt) = \nabla_{\bm \eta} I(W; \Yt) = \nabla_{\bm \eta} h(\Yt).
\label{eq:grad_mi_equals_grad_h}
\end{align}

To compute $\nabla_{\bm \eta} h(\Yt)$, we use the reparameterization $\Yt = f_{\bm \eta}(X) + Z_t$,
yielding 
\begin{align}
h(\Yt) = -\E_{X,Z_t}\left[\log p_{\Yt,\bm \eta}(f_{\bm \eta}(X) + Z_t)\right].
\end{align}

Under standard regularity conditions (differentiability of $f_{\bm \eta}$, 
exchangeability of differentiation and expectation), applying the chain rule yields
\begin{align}
\nabla_{\bm \eta} h(\Yt)
&= -\E_{X,Z_t}\left[\nabla_{\bm \eta} \log p_{\Yt,\bm \eta}(\bm y)\big|_{\bm y=f_{\bm \eta}(X)+Z_t}\right] \nonumber \\
&\quad - \E_{X,Z_t}\left[D f_{\bm \eta}(X)^\top \nabla_{\bm y} \log p_{\Yt,\bm \eta}(\bm y)\big|_{\bm y=f_{\bm \eta}(X)+Z_t}\right].
\label{eq:two_terms_proof}
\end{align}
The first term vanishes by the score property. Specifically, we have 
\begin{align}
\E_{Y_t}\left[\nabla_{\bm \eta} \log p_{\Yt,\bm \eta}(\Yt)\right]
&= \int \nabla_{\bm \eta} p_{\Yt,\bm \eta}(\bm y)\,d\bm y \\
&= \nabla_{\bm \eta} \int p_{\Yt,\bm \eta}(\bm y)\,d\bm y  \\
&= \nabla_{\bm \eta} 1 = \bm{0}.
\label{eq:score_property}
\end{align}
Substituting \eqref{eq:score_property} into \eqref{eq:two_terms_proof},
we have
\begin{align}
\nabla_{\bm \eta} h(\Yt) = -\E_{X,Z_t}\left[D f_{\bm \eta}(X)^\top \sYt(f_{\bm \eta}(X) + Z_t)\right].
\end{align}
\end{proof}

\end{document}
